\begin{proof}[Proof of \cref{obj:thm:minimax-lower}]
1. \emph{The three ingredients and the constant.}  By \cref{obj:lem:regret-separation}, there are constants \(c_{\mathrm{sep}}>0\) and \(N_1\) such that, for all \(n\ge N_1\) and every \(\pi\in\Pi\),
\[
\max_{\sigma\in\{+,-\}} R_{P_{n,\sigma}}(\pi)
\ge
c_{\mathrm{sep}}\,h_n^{1+\alpha}.
\]
By \cref{obj:lem:two-point-divergence}, there are \(C_\chi>0\) and \(N_2\) such that, for all \(n\ge N_2\), the product laws satisfy \(P_{n,+}^{\otimes n}\ll P_{n,-}^{\otimes n}\) with square-integrable density deviation and
\[
\chi^2(P_{n,+}^{\otimes n}\,\|\,P_{n,-}^{\otimes n})\le C_\chi .
\]
Applying \cref{obj:lem:le-cam-two-point-chisq} to this fixed budget gives a constant \(c_{\mathrm{test}}=c_{\mathrm{test}}(C_\chi)>0\) such that every measurable event \(A\) in the \(n\)-sample space satisfies
\[
P_{n,+}^{\otimes n}(A^c)+P_{n,-}^{\otimes n}(A)\ge c_{\mathrm{test}} .
\]
Set
\[
c=\frac{c_{\mathrm{sep}}\,c_{\mathrm{test}}}{2}>0 .
\]
% lean: minimax_lower

2. \emph{Membership and the rate identity.}  By \cref{obj:lem:witness-membership}, enlarging \(N_1,N_2\) if necessary, both witness laws \(P_{n,+}\) and \(P_{n,-}\) belong to \(\mathcal P_{\alpha,\gamma}\) for all \(n\) under consideration.  Since
\[
h_n=n^{-1/(2+\alpha+\beta_{\alpha,\gamma})},
\qquad
r_\star(\alpha,\gamma)=\frac{1+\alpha}{2+\alpha+\beta_{\alpha,\gamma}},
\]
raising \(h_n\) to the power \(1+\alpha\) gives the exact identity
\[
h_n^{1+\alpha}=n^{-r_\star(\alpha,\gamma)} .
\]
The policy class is nonempty, so the collection of measurable \(\Pi\)-valued estimators is nonempty (a constant estimator qualifies) and the infimum in \cref{obj:def:minimax-regret} is over a nonempty set; and regret is nonnegative and bounded by \(2\) under \cref{obj:ass:bounded-outcome}, so all displayed expectations are finite.  This is the only role of the nonemptiness and boundedness bookkeeping. % lean: witness_membership % lean: lawRegret_nonneg_of_wellformed % lean: lawRegret_le_two_of_wellformed

3. \emph{The threshold event.}  Fix such an \(n\), and let \(\widehat\pi\) be any measurable \(\Pi\)-valued estimator.  Write
\[
R_+
=
\mathbb E_{P_{n,+}^{\otimes n}} R_{P_{n,+}}(\widehat\pi),
\qquad
R_-
=
\mathbb E_{P_{n,-}^{\otimes n}} R_{P_{n,-}}(\widehat\pi),
\]
put
\[
s_n=c_{\mathrm{sep}}h_n^{1+\alpha}>0,
\qquad
A=\left\{\text{samples}:\ s_n\le R_{P_{n,-}}(\widehat\pi)\right\}.
\]
The set \(A\) is measurable because the map from the sample to \(R_{P_{n,-}}(\widehat\pi)\) is measurable, which is part of the estimator condition in \cref{obj:def:minimax-regret}. % lean: minimax_lower

4. \emph{Two Markov-type threshold bounds.}  For a probability measure \(Q\), a nonnegative \(Q\)-integrable function \(f\), and a level \(s>0\), the pointwise inequality \(f\ge s\,\mathbf 1\{f\ge s\}\) integrates to
\[
s\,Q\{f\ge s\}\le\int f\,dQ .
\]
Applying this to the nonnegative integrable regret maps under the two product laws,
\[
s_n\,P_{n,-}^{\otimes n}(A)\le R_-,
\qquad
s_n\,P_{n,+}^{\otimes n}\left\{s_n\le R_{P_{n,+}}(\widehat\pi)\right\}\le R_+ .
\]
% lean: mul_meas_ge_le_integral_of_nonneg
On \(A^c\) we have \(R_{P_{n,-}}(\widehat\pi)<s_n\); applying the separation bound of Step 1 pointwise to the realized policy \(\widehat\pi\in\Pi\), the maximum of the two regrets is at least \(s_n\), so it must be the plus-regret that attains it.  Hence
\[
A^c\subseteq \left\{s_n\le R_{P_{n,+}}(\widehat\pi)\right\},
\]
and by monotonicity of \(P_{n,+}^{\otimes n}\) together with the second threshold bound,
\[
s_n\,P_{n,+}^{\otimes n}(A^c)\le R_+ .
\]

5. \emph{Combining with the testing floor.}  Adding the two displayed bounds and inserting the testing lower bound of Step 1,
\[
s_n\,c_{\mathrm{test}}
\le
s_n\left\{P_{n,+}^{\otimes n}(A^c)+P_{n,-}^{\otimes n}(A)\right\}
\le
R_+ + R_- .
\]
Since \(R_++R_-\le 2\max\{R_+,R_-\}\), dividing by \(2\) gives
\[
\max\{R_+,R_-\}
\ge
\frac{s_n c_{\mathrm{test}}}{2}
=
\frac{c_{\mathrm{sep}}c_{\mathrm{test}}}{2}\,h_n^{1+\alpha}
=
c\,h_n^{1+\alpha}
=
c\,n^{-r_\star(\alpha,\gamma)} .
\]
% lean: le_cam_two_point_chisq

6. \emph{Passing to the minimax risk.}  Because \(P_{n,+},P_{n,-}\in\mathcal P_{\alpha,\gamma}\) by Step 2, each of \(R_+\) and \(R_-\) is one of the values over which the supremum in \cref{obj:def:minimax-regret} is taken, so
\[
\max\{R_+,R_-\}
\le
\sup_{P\in\mathcal P_{\alpha,\gamma}}
\mathbb E_{P^{\otimes n}} R_P(\widehat\pi).
\]
Therefore every admissible estimator satisfies
\[
\sup_{P\in\mathcal P_{\alpha,\gamma}}
\mathbb E_{P^{\otimes n}} R_P(\widehat\pi)
\ge
c\,n^{-r_\star(\alpha,\gamma)},
\]
with \(c\) not depending on \(\widehat\pi\) or on \(n\).  Taking the infimum over all measurable \(\Pi\)-valued estimators in \cref{obj:def:minimax-regret} gives
\[
M_n(\alpha,\gamma)\ge c\,n^{-r_\star(\alpha,\gamma)}
\]
for all sufficiently large \(n\). % lean: minimax_lower
\end{proof}
